\section{The Holy Ghost Movement}

Around the 11\textsuperscript{th} century, there was a manifestation of a new spirituality, attributed to the Holy Ghost. However, this led to commotion and agitation. Carl Jung describes it:

\begin{quotex}
Everyone felt the rushing wind of the pneuma; it was an age of new and unprecedented ideas which were blazoned abroad by the Cathari, Patarenes, Concorricci, Waldenses, Poor Men of Lyons, Beghards, Brethren of the Free Spirit, ``Bread through God,'' and whatever else these movements were called.
\flright{\textsc{Carl Jung}, \emph{Aion}}
\end{quotex}


There beliefs can be summarized:

\begin{itemize}


\item they believe themselves to be God by nature without distinction

\item they are eternal.

\item they have no need of God or the Godhead.

\item they constitute the kingdom of heaven.

\item they are immutable in the new rock, that they rejoice in naught and are troubled by naught.

\item a man is bound to follow his inner instinct rather than the truth of the Gospel which is preached every day

\item they believe the Gospel to contain poetical matters which are not true.
\end{itemize}


In summary:

\begin{quotex}
They were made up of people who identified themselves (or were identified) with God, who deemed themselves supermen, had a critical approach to the gospels, followed the promptings of the inner man, and understood the kingdom of heaven to be within. In a sense, therefore, they were modern in their outlook, but they had a religious inflation instead of the rationalistic and political psychosis that is the affliction of our day.
\end{quotex}


In that milieu, \textbf{Joachim of Fiore} announced the coming of the new Age of the Spirit, following the preceding Age of the Son and the Age of the Father. This third age, he claimed, began with Saint Benedict. Although his theology was ambiguous, his life was pious. Dante even recognized him in Paradise. Nevertheless, the effects were real, as described by Jung. The Age of the Spirit is upon us, as well as the seductions of the Antichrist. Deep discernment is necessary to tell the difference. So he cannot be held responsible for the excesses of that movement. It took a while for the best theologians to digest the Holy Ghost Movement.

\begin{quotex}
The repercussions of the Holy Ghost movement spread, in the years that followed, to four minds of immense significance for the future. These were Albertus Magnus (1193-1280); his pupil Thomas Aquinas, the philosopher of the Church and an adept in alchemy (as also was Albertus); Roger Bacon (c. 1214-c.1294), the English forerunner of inductive science; and finally Meister Eckhart (c. 1260-1327), the independent religious thinker, now enjoying a real revival after six hundred years of obscurity.
\end{quotex}


\paragraph{Levels of Being}


\begin{quotex}
Man is as if he were three men, an \textbf{animal man}, as he is according to the senses, then a \textbf{rational man}, and finally the \textbf{highest god-like man}. . . One is the external, animal, sensual man; the other is the internal, rational man, with his rational faculties; the third man is the spirit, the highest part of the soul.
\flright{\textsc{Johannes Tauler}}
\end{quotex}


Although Tauler's remark is consistent with St. Paul's teaching, it is never heard in exoteric preaching. Although in most cases this teaching is either unknown or poorly understood, there are sound reasons for suppressing it. Obviously, it is incompatible with the egalitarian attitude so prevalent today.

The main problem is that people will tend to overestimate their own progress. Anyone can claim, for example, to be the highest type of man, apart from any actual attainment. The Cathars used it to create a privileged class. Although people of the spiritual class claimed to transcend personal desires, in actuality they were libertines. This was justified by the claim that such liberties had no effect on them. The greatest danger is that the ``animal man'' may be associated with a particular race or ethnicity, i.e., a biological rather than spiritual state.

We had a recent comment from a reader who asserted that his spiritual director advised him not to read the mystics directly. This demonstrates the point about suppression.

There are Greek terms---hylic, psychical, pneumatic---for Tauler's classification, although ``carnal'' is often used for ``hylic''. The effect of religious teachings is different when they affect the different types. For example,

\begin{itemize}


\item \textbf{Hylic}. The hylic personality will be most impressed with material and physical issues. For example, he or she will be interested in miracles or other unusual phenomena, patterns of clouds in the sky, historical events (which are for some reason called ``literal''), and so on.

\item \textbf{Rational}. This type loves to argue. He misunderstands spiritual combat to be some sort of intellectual battle of ideas. Instead of seeking wisdom like the philosophers, he enjoys ``making a case'', as the sophists would do. He maintains centuries old divisions, regarding them as inspired somehow, instead of looking for reconciliation.

\item \textbf{Spiritual}. The spiritual man wonders whether he should laugh or cry.
\end{itemize}


\paragraph{Psychic Reality}


Like \textbf{Meister Eckhart}, Tauler understands that the real birth of the Son is in the soul when he writes:

\begin{quotex}
God accomplishes all His works in the soul and gives them to the soul; and the Father brings forth His only-begotten Son in the soul, as truly as He brings Him forth in eternity, neither less, nor more.
\end{quotex}


That would disturb the materially minded, the hylic personalities, who believe that the physical is more real than the psychical. Whatever his faults and limitations may be, Carl Jung seems closer to Eckhart and Tauler in his insistence that the psyche is \emph{real}. For many, the notion that something happens in the soul means that it is less real, it is imaginary, it is all ``in your head''. In opposition to that notion, Jung makes this claim:

\begin{quotex}
It does not matter at all that a physically impossible fact is asserted, because all religious assertions are physical impossibilities. If they were not so, they would necessarily be treated in textbooks of natural science. But religious statements without exception have to do with the reality of the psyche and not with the reality of \emph{physis}.
\end{quotex}


Now a \textbf{Rene Guenon} might object that religious statements are disguised metaphysical statements. Nevertheless, Jung's observation is a good half way point, and bears more fruit in practice.

\paragraph{Science and the Carnal Man}


Rudolf Steiner, in his commentary on Tauler, describes the limitations of science:

\begin{quotex}
Man is entangled in the world of the senses and in the laws of nature, by which the world of the senses is dominated. He himself is a result of this world. He lives because its forces and substances are active in him, and he perceives and judges this world of the senses in accordance with the laws by which it and he are constructed.
\end{quotex}


This is the circularity of the deracinated, secularized, rational, scientific man. He is fundamentally a carnal man, a highly intelligent animal. However, all his science is self-confirming, so he is never able to break out of the circle to reach transcendent, or higher, knowledge. Steiner notes this anomaly:

\begin{quotex}
Do we not stand above all mere conformity to natural laws when we survey how we ourselves are integrated into nature? We \emph{see} with our eye in accordance with the laws of nature. But we also \emph{understand} the laws in accordance with which we see. We can stand on a higher elevation and survey simultaneously the external world and ourselves in interplay. Is not then a nature active within us which is higher than the sensory-organic personality which acts according to natural laws and with natural laws? In such activity is there still a partition between our inner world and the external world?
\end{quotex}


A recent example is this study that \textit{Finds It Might Not Be Consciousness That Drives The Human Mind}\footnote{\url{https://www.sciencealert.com/research-finds-it-might-not-be-consciousness-that-drives-the-human-mind}}. The research notes that the carnal mind ``gives us a feeling of ownership and control over the thoughts, emotions and experiences that we have every day''. On the contrary, the esoterist, the Hermetist, the yogi have known that for millennia, so the scientists are wasting their time. The mystics understand that thoughts and feelings spontaneously arise, in a manner not unlike our experiences of the external world, so that they are not part of our ``inner being''. They have developed practices to observe the arising of thought and the means to resist the internal programming that they represent.

Their conclusion is quite odd:

\begin{quotex}
the contents of consciousness are generated ``behind the scenes'' by fast, efficient, non-conscious systems in our brains.
\end{quotex}


We are left wondering how exactly does the brain produce thoughts. Is there some chemical formula or equation of physics that describes each thought? Since our lived world is the creation of our thoughts, it would be incredible if someone could describe a modern city as the result of some efficient system in the brain.

The denial of the immateriality of the soul is a misuse of intelligence.


\flrightit{Posted on 2022-05-01 by Cologero}

\begin{center}* * *\end{center}

\begin{footnotesize}
\begin{sffamily}
\texttt{De Bosit on 2022-05-02 at 00:55 said:}

I like the studies conclusion.

I think of it as frequencies, where you let the higher regions do their thing in peace and I just say what comes to mind without deliberation.

If you make those higher regions known to yourself you will come off like coked up Charley Sheen, well at least that is my experience.

Better let it do it's thing in peace and enjoy the comfort.

\hfill


\end{sffamily}
\end{footnotesize}

